\documentclass[a4paper,12pt,twoside]{memoir}

% Castellano
\usepackage[spanish,es-tabla]{babel}
\selectlanguage{spanish}
\usepackage[utf8]{inputenc}
\usepackage[T1]{fontenc}
\usepackage{lmodern} % Scalable font
\usepackage{microtype}
\usepackage{placeins}
\usepackage{floatrow}

\RequirePackage{booktabs}
\RequirePackage[table]{xcolor}
\RequirePackage{xtab}
\RequirePackage{multirow}

% Links
\PassOptionsToPackage{hyphens}{url}\usepackage[colorlinks]{hyperref}
\hypersetup{
	allcolors = {red}
}

% Ecuaciones
\usepackage{amsmath}

% Rutas de fichero / paquete
\newcommand{\ruta}[1]{{\sffamily #1}}

% Párrafos
\nonzeroparskip

% Huérfanas y viudas
\widowpenalty100000
\clubpenalty100000

% Imágenes

% Comando para insertar una imagen en un lugar concreto.
% Los parámetros son:
% 1 --> Ruta absoluta/relativa de la figura
% 2 --> Texto a pie de figura
% 3 --> Tamaño en tanto por uno relativo al ancho de página
\usepackage{graphicx}
\newcommand{\imagen}[3]{
	\begin{figure}[!h]
		\centering
		\includegraphics[width=#3\textwidth]{#1}
		\caption{#2}\label{fig:#1}
	\end{figure}
	\FloatBarrier
}

% Comando para insertar una imagen sin posición.
% Los parámetros son:
% 1 --> Ruta absoluta/relativa de la figura
% 2 --> Texto a pie de figura
% 3 --> Tamaño en tanto por uno relativo al ancho de página
\newcommand{\imagenflotante}[3]{
	\begin{figure}
		\centering
		\includegraphics[width=#3\textwidth]{#1}
		\caption{#2}\label{fig:#1}
	\end{figure}
}

% El comando \figura nos permite insertar figuras comodamente, y utilizando
% siempre el mismo formato. Los parametros son:
% 1 --> Porcentaje del ancho de página que ocupará la figura (de 0 a 1)
% 2 --> Fichero de la imagen
% 3 --> Texto a pie de imagen
% 4 --> Etiqueta (label) para referencias
% 5 --> Opciones que queramos pasarle al \includegraphics
% 6 --> Opciones de posicionamiento a pasarle a \begin{figure}
\newcommand{\figuraConPosicion}[6]{%
  \setlength{\anchoFloat}{#1\textwidth}%
  \addtolength{\anchoFloat}{-4\fboxsep}%
  \setlength{\anchoFigura}{\anchoFloat}%
  \begin{figure}[#6]
    \begin{center}%
      \Ovalbox{%
        \begin{minipage}{\anchoFloat}%
          \begin{center}%
            \includegraphics[width=\anchoFigura,#5]{#2}%
            \caption{#3}%
            \label{#4}%
          \end{center}%
        \end{minipage}
      }%
    \end{center}%
  \end{figure}%
}

%
% Comando para incluir imágenes en formato apaisado (sin marco).
\newcommand{\figuraApaisadaSinMarco}[5]{%
  \begin{figure}%
    \begin{center}%
    \includegraphics[angle=90,height=#1\textheight,#5]{#2}%
    \caption{#3}%
    \label{#4}%
    \end{center}%
  \end{figure}%
}
% Para las tablas
\newcommand{\otoprule}{\midrule [\heavyrulewidth]}
%
% Nuevo comando para tablas pequeñas (menos de una página).
\newcommand{\tablaSmall}[5]{%
 \begin{table}
  \begin{center}
   \rowcolors {2}{gray!35}{}
   \begin{tabular}{#2}
    \toprule
    #4
    \otoprule
    #5
    \bottomrule
   \end{tabular}
   \caption{#1}
   \label{tabla:#3}
  \end{center}
 \end{table}
}

%
% Nuevo comando para tablas pequeñas (menos de una página).
\newcommand{\tablaSmallSinColores}[5]{%
 \begin{table}[H]
  \begin{center}
   \begin{tabular}{#2}
    \toprule
    #4
    \otoprule
    #5
    \bottomrule
   \end{tabular}
   \caption{#1}
   \label{tabla:#3}
  \end{center}
 \end{table}
}

\newcommand{\tablaApaisadaSmall}[5]{%
\begin{landscape}
  \begin{table}
   \begin{center}
    \rowcolors {2}{gray!35}{}
    \begin{tabular}{#2}
     \toprule
     #4
     \otoprule
     #5
     \bottomrule
    \end{tabular}
    \caption{#1}
    \label{tabla:#3}
   \end{center}
  \end{table}
\end{landscape}
}

%
% Nuevo comando para tablas grandes con cabecera y filas alternas coloreadas en gris.
\newcommand{\tabla}[6]{%
  \begin{center}
    \tablefirsthead{
      \toprule
      #5
      \otoprule
    }
    \tablehead{
      \multicolumn{#3}{l}{\small\sl continúa desde la página anterior}\\
      \toprule
      #5
      \otoprule
    }
    \tabletail{
      \hline
      \multicolumn{#3}{r}{\small\sl continúa en la página siguiente}\\
    }
    \tablelasttail{
      \hline
    }
    \bottomcaption{#1}
    \rowcolors {2}{gray!35}{}
    \begin{xtabular}{#2}
      #6
      \bottomrule
    \end{xtabular}
    \label{tabla:#4}
  \end{center}
}

%
% Nuevo comando para tablas grandes con cabecera.
\newcommand{\tablaSinColores}[6]{%
  \begin{center}
    \tablefirsthead{
      \toprule
      #5
      \otoprule
    }
    \tablehead{
      \multicolumn{#3}{l}{\small\sl continúa desde la página anterior}\\
      \toprule
      #5
      \otoprule
    }
    \tabletail{
      \hline
      \multicolumn{#3}{r}{\small\sl continúa en la página siguiente}\\
    }
    \tablelasttail{
      \hline
    }
    \bottomcaption{#1}
    \begin{xtabular}{#2}
      #6
      \bottomrule
    \end{xtabular}
    \label{tabla:#4}
  \end{center}
}

%
% Nuevo comando para tablas grandes sin cabecera.
\newcommand{\tablaSinCabecera}[5]{%
  \begin{center}
    \tablefirsthead{
      \toprule
    }
    \tablehead{
      \multicolumn{#3}{l}{\small\sl continúa desde la página anterior}\\
      \hline
    }
    \tabletail{
      \hline
      \multicolumn{#3}{r}{\small\sl continúa en la página siguiente}\\
    }
    \tablelasttail{
      \hline
    }
    \bottomcaption{#1}
  \begin{xtabular}{#2}
    #5
   \bottomrule
  \end{xtabular}
  \label{tabla:#4}
  \end{center}
}



\definecolor{cgoLight}{HTML}{EEEEEE}
\definecolor{cgoExtralight}{HTML}{FFFFFF}

%
% Nuevo comando para tablas grandes sin cabecera.
\newcommand{\tablaSinCabeceraConBandas}[5]{%
  \begin{center}
    \tablefirsthead{
      \toprule
    }
    \tablehead{
      \multicolumn{#3}{l}{\small\sl continúa desde la página anterior}\\
      \hline
    }
    \tabletail{
      \hline
      \multicolumn{#3}{r}{\small\sl continúa en la página siguiente}\\
    }
    \tablelasttail{
      \hline
    }
    \bottomcaption{#1}
    \rowcolors[]{1}{cgoExtralight}{cgoLight}

  \begin{xtabular}{#2}
    #5
   \bottomrule
  \end{xtabular}
  \label{tabla:#4}
  \end{center}
}



\graphicspath{ {./img/} }

% Capítulos
\chapterstyle{bianchi}
\newcommand{\capitulo}[2]{
	\setcounter{chapter}{#1}
	\setcounter{section}{0}
	\setcounter{figure}{0}
	\setcounter{table}{0}
	\chapter*{\thechapter.\enskip #2}
	\addcontentsline{toc}{chapter}{\thechapter.\enskip #2}
	\markboth{#2}{#2}
}

% Apéndices
\renewcommand{\appendixname}{Apéndice}
\renewcommand*\cftappendixname{\appendixname}

\newcommand{\apendice}[1]{
	%\renewcommand{\thechapter}{A}
	\chapter{#1}
}

\renewcommand*\cftappendixname{\appendixname\ }

% Formato de portada
\makeatletter
\usepackage{xcolor}
\newcommand{\tutor}[1]{\def\@tutor{#1}}
\newcommand{\course}[1]{\def\@course{#1}}
\definecolor{cpardoBox}{HTML}{E6E6FF}
\def\maketitle{
  \null
  \thispagestyle{empty}
  % Cabecera ----------------
\noindent\includegraphics[width=\textwidth]{cabecera}\vspace{1cm}%
  \vfill
  % Título proyecto y escudo informática ----------------
  \colorbox{cpardoBox}{%
    \begin{minipage}{.8\textwidth}
      \vspace{.5cm}\Large
      \begin{center}
      \textbf{TFG del Grado en Ingeniería Informática}\vspace{.6cm}\\
      \textbf{\LARGE\@title{}}
      \end{center}
      \vspace{.2cm}
    \end{minipage}

  }%
  \hfill\begin{minipage}{.20\textwidth}
    \includegraphics[width=\textwidth]{escudoInfor}
  \end{minipage}
  \vfill
  % Datos de alumno, curso y tutores ------------------
  \begin{center}%
  {%
    \noindent\LARGE
    Presentado por \@author{}\\ 
    en Universidad de Burgos --- \@date{}\\
    Tutor: \@tutor{}\\
  }%
  \end{center}%
  \null
  \cleardoublepage
  }
\makeatother

\newcommand{\nombre}{Javier Pampliega García} %%% cambio de comando

% Datos de portada
\title{PlanLMS \\ 
        Aplicación multiplataforma para la gestión de tareas de la plataforma Moodle}
\author{\nombre}
\tutor{Raúl Marticorena Sánchez}
\date{\today}

\begin{document}

\maketitle


\newpage\null\thispagestyle{empty}\newpage


%%%%%%%%%%%%%%%%%%%%%%%%%%%%%%%%%%%%%%%%%%%%%%%%%%%%%%%%%%%%%%%%%%%%%%%%%%%%%%%%%%%%%%%%
\thispagestyle{empty}


\noindent\includegraphics[width=\textwidth]{cabecera}\vspace{1cm}

\noindent D. Raúl Marticorena Sánchez, profesor del departamento de Ingeniería Informática, área de Lenguajes y Sistemas Informáticos.

\noindent Expone:

\noindent Que el alumno D. \nombre, con DNI 71755086P, ha realizado el Trabajo final de Grado en Ingeniería Informática titulado PlanLMS Aplicación multiplataforma para la gestión de tareas de la plataforma Moodle. 

\noindent Y que dicho trabajo ha sido realizado por el alumno bajo la dirección del que suscribe, en virtud de lo cual se autoriza su presentación y defensa.

\begin{center} %\large
En Burgos, {\large \today}
\end{center}

\vfill\vfill\vfill

% Author and supervisor
%\begin{minipage}{0.45\textwidth}
%\begin{flushleft} %\large
%Vº. Bº. del Tutor:\\[2cm]
%D. nombre tutor
%\end{flushleft}
%\end{minipage}
%\hfill
%\begin{minipage}{0.45\textwidth}
%\begin{flushleft} %\large
%Vº. Bº. del co-tutor:\\[2cm]
%D. nombre co-tutor
%\end{flushleft}
%\end{minipage}
%\hfill

%\vfill

% para casos con solo un tutor comentar lo anterior
% y descomentar lo siguiente
\begin{center}
Vº. Bº. del Tutor:\\[2cm]
D. Raúl Marticorena Sánchez
\end{center}


\newpage\null\thispagestyle{empty}\newpage




\frontmatter

% Abstract en castellano
\renewcommand*\abstractname{Resumen}
\begin{abstract}
PlanLMS es una aplicación multiplataforma desarrollada con Flutter con el objetivo de facilitar la planificación de tareas a los usuarios de plataformas Moodle. Su funcionamiento es posible por la utilización de los servicios web de Moodle. 

A través de una interfaz simple e intuitiva, se busca que el usuario logre planificar sus actividades por medio de un diagrama de Gantt que le permita visualizar la distribución de estas actividades en el tiempo. Por otra parte, incorpora una sección donde el usuario puede visualizar todas sus actividades de forma ordenada. Además, se incorpora otra funcionalidad que permite al usuario crear sus propias tareas asociadas a un curso.

Todas estas funcionalidades pueden ser reguladas mediante el sistema de filtrado que permite que el usuario visualice únicamente el contenido que necesite.


\end{abstract}

\renewcommand*\abstractname{Descriptores}
\begin{abstract}
Gestión de tareas, Moodle, diagrama de Gantt, multiplataforma, Flutter, Dart.
\end{abstract}

\clearpage

% Abstract en inglés
\renewcommand*\abstractname{Abstract}
\begin{abstract}
PlanLMS is a cross-platform application developed with Flutter, designed to make task planning easier for Moodle users. It operates by using Moodle´s web services.

Through a simple and intuitive interface, it allows the user to plan activities via a Gantt chart that visualizes their distribution over time. Likewise, it includes a section where the user can view all their activities in an organized list. Additionally, it offers a feature that lets the user create their own tasks linked to a course.

All of these functions can be managed through a filtering system that enables users to view only the content they need.
\end{abstract}

\renewcommand*\abstractname{Keywords}
\begin{abstract}
Task management, Moodle, Gantt chart, cross-platform, Flutter, Dart.
\end{abstract}

\clearpage

% Indices
\tableofcontents

\clearpage

\listoffigures

\clearpage

\listoftables
\clearpage

\mainmatter
\capitulo{1}{Introducción}

Durante los últimos años la implementación de tecnologías en el ámbito educativo ha transformado los procedimientos de enseñanza. Las plataformas de gestión de aprendizaje (Learning Management Systems o LMS) como Moodle se han consolidado como herramientas imprescindibles en la enseñanza y en el seguimiento del progreso académico, ya que facilitan la interacción entre alumnos y docentes, organización del contenido académico, evaluación de actividades, etc. Aunque el uso de Moodle es cómodo para el alumnado, cabe recalcar que carece de algunas funcionalidades que pueden enriquecer el uso de la plataforma.

La motivación que impulsa el desarrollo de este proyecto viene de mejorar la gestión de las actividades de Moodle, para dar una visión más clara al alumnado de las tareas que debe realizar, previniendo que las labores pendientes queden desatendidas. Además, la aplicación incorpora un sistema de filtrado que proporciona un abanico de amplias posibilidades en el momento de filtrar el contenido que el usuario desee.

Por ello, en este Trabajo de Fin de Grado se propone la implementación de una aplicación multiplataforma desarrollada con Dart y Flutter, con el objetivo de complementar y mejorar la experiencia que ofrece Moodle. Para lograr esto último es necesario recurrir a los servicios web ofrecidos por Moodle, lo que permite la autenticación de usuarios, consulta de cursos, actividades a realizar o calificaciones obtenidas, entre decenas de funciones. Por otro lado, se presta especial atención al diseño de la interfaz de usuario, así como a la experiencia de usuario, mejorando la presentación y el uso de la aplicación.

\subsection{Estructura del material entregado}
A continuación se detallarán las estructuras de los materiales entregados. Desde la memoria y anexos, hasta la estructura de directorios y ficheros alojados en el repositorio de GitHub.

\subsubsection{Estructura de la memoria}

\begin{itemize}
    \item \textbf{Introducción:} proporciona una visión general del contexto del proyecto, además de la estructura que lo compone.
    \item \textbf{Objetivos del proyecto:} describe los objetivos que se persiguen para alcanzar las metas establecidas.
    \item \textbf{Conceptos teóricos:} explica conceptos teóricos necesarios para la compresión total del proyecto y los objetivos que persigue.
    \item \textbf{Técnicas y herramientas:} describe las técnicas y herramientas empleadas para lograr el desarrollo del proyecto.
    \item \textbf{Aspectos relevantes del desarrollo del proyecto:} detalla los aspectos más destacados durante el desarrollo del proyecto, desde los inconvenientes hasta las decisiones tomadas para resolver los desafíos del proyecto.
    \item \textbf{Trabajos relacionados:} revisión de proyectos similares, realizando comparativas con dichos proyectos.
    \item \textbf{Conclusiones y líneas de trabajo futuras:} recoge las conclusiones finales del proyecto y especifica posibles mejoras a implementar en el futuro.
\end{itemize}

\subsubsection{Estructura de los anexos}

\begin{itemize}
    \item \textbf{Plan de proyecto software:} detalla la planificación del proyecto y estudia las viabilidades económicas y legales del mismo.
    \item \textbf{Especificación de requisitos:} detalla los requisitos y casos de uso del proyecto.
    \item \textbf{Especificación de diseño:} especifica los aspectos de diseño, detallando la interfaz de usuario, diseño de datos, diagrama de clases, etc.
    \item \textbf{Documentación técnica de programación:} aporta documentación con la explicación de aspectos técnicos del proyecto para programadores.
    \item \textbf{Documentación de usuario:} aporta documentación relevante para que el usuario tenga conocimiento sobre el uso de la aplicación.
\end{itemize}

\capitulo{2}{Objetivos del proyecto}

El objetivo principal del proyecto es el desarrollo de una aplicación multiplataforma con Dart y Flutter, que permita al alumnado de distintas instituciones educativas planificar las actividades de su plataforma Moodle. Se busca potenciar la organización del usuario para incrementar su rendimiento académico, así como sus capacidades de planificación de tareas.

El proyecto se centra en uso de un diagrama de Gantt que facilita la visualización temporal de las actividades a realizar. Por otra parte, se implementa la creación de tareas con niveles de prioridad y ligadas a un curso matriculado. Estas funcionalidades pueden ser filtradas al gusto del usuario con el sistema de filtrado incorporado.

\begin{itemize}
    \item \textbf{Visualización gráfica de actividades:} El objetivo principal es la implementación de un diagrama de Gantt interactivo que permita al usuario visualizar de forma clara y ordenada las fechas y duración de las tareas, de las actividades de Moodle. Además, el usuario podrá desplazarse a través del diagrama y podrá ver sus actividades ordenadas por cursos en un sección pareja al diagrama.
    \item \textbf{Creación de tareas propias:} La aplicación ofrecerá al usuario la posibilidad de definir sus propias tareas ligadas a un curso en el que se encuentre matriculado y podrá asignarles unas prioridad entre otras opciones. Todas las tareas creadas por parte del usuario se mostrarán por pantalla en función de su estado: Pendiente o Completada. Además, estarán ordenadas en función de su fecha de caducidad o su fecha de finalización.
    \item \textbf{Filtrado general de actividades y tareas:} Una de las funcionalidades más importantes será el filtrado de todos los eventos de ambas funcionalidades mencionadas anteriormente. Esto será posible gracias al sistema de filtrado disponible en la pantalla inicial. Desde este, el usuario podrá filtrar por cursos, fechas, tipos de actividad y disponibilidad de fechas de apertura y cierre. Con ello se pretende que el usuario modifique el contenido de su aplicación a gusto propio y así logre mejorar su capacidad de organización.
\end{itemize}

\subsection{Objetivos personales}
En esta sección se establecen los objetivos personales que se pretenden superar en la realización de este proyecto:

\begin{itemize}
    \item Aprendizaje en el desarrollo de aplicaciones, en este caso multiplataforma, para futuros proyectos.
    \item Mejorar el manejo de APIs.
    \item Entender el uso de la metodología SCRUM para su uso aplicado en proyectos reales.
    \item Mejorar la lógica de programación para desarrollos complejos.
    \item Descubrir y explorar nuevas herramientas que no había empleado con anterioridad.
\end{itemize}

\subsection{Objetivos técnicos}
En esta sección se especifican los objetivos técnicos que se pretenden alcanzar con la realización del proyecto:

\begin{itemize}
    \item Uso de GitHub con GitBash para el control de versiones.
    \item Uso de Visual Studio Code para el desarrollo del proyecto.
    \item Uso de Dart y Flutter para el desarrollo de la aplicación multiplataforma.
    \item Uso de Postman para realizar peticiones a la API.
    \item Uso de AndroidEmulator para la simulación de un entorno Android.
    \item Uso de Figma para el prototipo de la aplicación.
    \item Uso de Overleaf para la realización de la documentación.
\end{itemize}
\capitulo{3}{Conceptos teóricos}

En aquellos proyectos que necesiten para su comprensión y desarrollo de unos conceptos teóricos de una determinada materia o de un determinado dominio de conocimiento, debe existir un apartado que sintetice dichos conceptos.

Algunos conceptos teóricos de \LaTeX{} \footnote{Créditos a los proyectos de Álvaro López Cantero: Configurador de Presupuestos y Roberto Izquierdo Amo: PLQuiz}.

\section{Secciones}

Las secciones se incluyen con el comando section.

\subsection{Subsecciones}

Además de secciones tenemos subsecciones.

\subsubsection{Subsubsecciones}

Y subsecciones. 


\section{Referencias}

Las referencias se incluyen en el texto usando cite~\cite{wiki:latex}. Para citar webs, artículos o libros~\cite{koza92}, si se desean citar más de uno en el mismo lugar~\cite{bortolot2005, koza92}.


\section{Imágenes}

Se pueden incluir imágenes con los comandos standard de \LaTeX, pero esta plantilla dispone de comandos propios como por ejemplo el siguiente:

\imagen{escudoInfor}{Autómata para una expresión vacía}{.5}



\section{Listas de items}

Existen tres posibilidades:

\begin{itemize}
	\item primer item.
	\item segundo item.
\end{itemize}

\begin{enumerate}
	\item primer item.
	\item segundo item.
\end{enumerate}

\begin{description}
	\item[Primer item] más información sobre el primer item.
	\item[Segundo item] más información sobre el segundo item.
\end{description}
	
\begin{itemize}
\item 
\end{itemize}

\section{Tablas}

Igualmente se pueden usar los comandos específicos de \LaTeX o bien usar alguno de los comandos de la plantilla.

\tablaSmall{Herramientas y tecnologías utilizadas en cada parte del proyecto}{l c c c c}{herramientasportipodeuso}
{ \multicolumn{1}{l}{Herramientas} & App AngularJS & API REST & BD & Memoria \\}{ 
HTML5 & X & & &\\
CSS3 & X & & &\\
BOOTSTRAP & X & & &\\
JavaScript & X & & &\\
AngularJS & X & & &\\
Bower & X & & &\\
PHP & & X & &\\
Karma + Jasmine & X & & &\\
Slim framework & & X & &\\
Idiorm & & X & &\\
Composer & & X & &\\
JSON & X & X & &\\
PhpStorm & X & X & &\\
MySQL & & & X &\\
PhpMyAdmin & & & X &\\
Git + BitBucket & X & X & X & X\\
Mik\TeX{} & & & & X\\
\TeX{}Maker & & & & X\\
Astah & & & & X\\
Balsamiq Mockups & X & & &\\
VersionOne & X & X & X & X\\
} 

\capitulo{4}{Técnicas y herramientas}

Esta parte de la memoria tiene como objetivo presentar las técnicas metodológicas y las herramientas de desarrollo que se han utilizado para llevar a cabo el proyecto. Si se han estudiado diferentes alternativas de metodologías, herramientas, bibliotecas se puede hacer un resumen de los aspectos más destacados de cada alternativa, incluyendo comparativas entre las distintas opciones y una justificación de las elecciones realizadas. 
No se pretende que este apartado se convierta en un capítulo de un libro dedicado a cada una de las alternativas, sino comentar los aspectos más destacados de cada opción, con un repaso somero a los fundamentos esenciales y referencias bibliográficas para que el lector pueda ampliar su conocimiento sobre el tema.



\capitulo{5}{Aspectos relevantes del desarrollo del proyecto}

Este apartado pretende recoger los aspectos más interesantes del desarrollo del proyecto, comentados por los autores del mismo.
Debe incluir desde la exposición del ciclo de vida utilizado, hasta los detalles de mayor relevancia de las fases de análisis, diseño e implementación.
Se busca que no sea una mera operación de copiar y pegar diagramas y extractos del código fuente, sino que realmente se justifiquen los caminos de solución que se han tomado, especialmente aquellos que no sean triviales.
Puede ser el lugar más adecuado para documentar los aspectos más interesantes del diseño y de la implementación, con un mayor hincapié en aspectos tales como el tipo de arquitectura elegido, los índices de las tablas de la base de datos, normalización y desnormalización, distribución en ficheros3, reglas de negocio dentro de las bases de datos (EDVHV GH GDWRV DFWLYDV), aspectos de desarrollo relacionados con el WWW...
Este apartado, debe convertirse en el resumen de la experiencia práctica del proyecto, y por sí mismo justifica que la memoria se convierta en un documento útil, fuente de referencia para los autores, los tutores y futuros alumnos.

\capitulo{6}{Trabajos relacionados}

En este apartado se presentan distintos trabajos o proyectos similares a la aplicación desarrollada. El objetivo es realizar una breve comparativa de las opciones que ofrecen distintas alternativas en comparación a la desarrollada, con la finalidad de dar una visión al usuario de las opciones que puede utilizar en función de sus necesidades.

\section{Aplicación móvil de Moodle}
La aplicación oficial de Moodle ofrece todas las funcionalidades que ofrece la versión web, pero adaptada a dispositivos móviles \cite{related_moodle}. La aplicación de Moodle incorpora las siguientes funciones:
\begin{itemize}
    \item Acceso a los cursos.
    \item Comunicación con otros participantes de los cursos.
    \item Envío de tareas.
    \item Seguimiento del progreso.
    \item Realización de actividades.
\end{itemize}

\section{Cuckoo}
Se trata de un proyecto desarrollado en Flutter cuya finalidad es ofrecer un rendimiento mejorado en comparación a la aplicación oficial de Moodle. Además, implementa funcionalidades que pueden resultar muy útiles, así como puede ser modificado para dar soporte a otras plataformas de Moodle \cite{related_cuckoo}.

Algunas funcionalidades que ofrece son:
\begin{itemize}
    \item \textbf{Autenticación:} hace uso del mismo sistema de autenticación de Moodle con el fin de garantizar la seguridad de la información de los usuarios.
    \item \textbf{Estado de finalización:} Cuckoo se sincroniza con Moodle para comprobar si los eventos con fecha limite están completados y, según corresponda, marcarlas en las lista de eventos de la propia aplicación.
    \item \textbf{Eventos personalizados:} los usuarios pueden crear sus propios eventos y vincularlos con un curso en el que se encuentre inscrito.
    \item \textbf{Acceso a los recursos de Moodle:} dispone de acceso a los contenidos descargables de los cursos, y permite descargarlos sin necesidad de acceder a Moodle.
\end{itemize}

\section{Trello}
Trello es un software de gestión en linea, que hace uso de la metodología Kanban. Esta metodología se basa en el uso de tarjetas de trabajo en un tablero, similar a una línea de producción de tareas con un \textit{status} \cite{related_trello}.

Trello sirve para organizar y gestionar tareas, ya sean proyectos o tareas del día, entre otros. Trello dispone de una gran cantidad de características que hacen de su uso una comodidad:
\begin{itemize}
    \item Sistema de trabajo colaborativo.
    \item Asignación de tareas.
    \item Organizar tareas por etiquetas.
    \item Multiplataforma.
    \item Fechas de vencimiento de tareas.
    \item Personalización de tableros.
\end{itemize}

\section{Tabla de comparativas}
A continuación se presenta una tabla comparando distintas funcionalidades de las aplicaciones presentadas anteriormente. El cometido de esta tabla es realizar una comparativa de los puntos fuertes de este proyecto, comparando las posibilidades del resto.

\tablaSmallSinColores
{Comparativa de puntos fuertes del proyecto con otros proyectos}
{ccccc}
{comparativa-proyectos}
{
    \textbf{Funcionalidad} \textbackslash \textbf{App} & \textbf{Moodle} & \textbf{Cuckoo} & \textbf{Trello} & \textbf{Proyecto} \\
}
{
    Autenticación & X & X & X & X \\
    Conexión Moodle & X & X & - & X \\
    Diagrama Gantt & - & - & - & X \\
    Filtro General & - & - & - & X \\
    Tareas Personales & X & X & X & X \\
}
\capitulo{7}{Conclusiones y Líneas de trabajo futuras}

\section{Conclusiones}
Una vez finalizado el desarrollo del proyecto de la aplicación multiplataforma para gestionar tareas de Moodle, puedo afirmar que se han cumplido con todos los objetivos propuestos y, además, se han incorporado funcionalidades que no estaban previstas desde comienzo del proyecto. Por lo tanto, se han completado con todos los objetivos del proyecto. El hecho de haber alcanzado las expectativas previstas supone un logro gratificante, que es fruto de una planificación en la que se realizaban constantes revisiones así como nuevas propuestas.

En el aspecto personal, me siento satisfecho con el resultado final, ya que se ofrece una aplicación al alumnado que puede usar de forma que le sirva para la planificación de su tiempo en el momento de realizar sus tareas. El hecho de haber resuelto todos los obstáculos presentes durante la realización del proyecto, ha supuesto un logro personal, donde me he demostrado a mi mismo que con constancia todo es posible.

Por otra parte, ha sido positivo adquirir nuevos conocimientos relacionados con el desarrollo de aplicaciones multiplataforma. Además, me ha motivado a seguir ampliando mis conocimientos acerca de esta tecnología. Estos conocimientos podrán ser utilizados en el futuro para el desarrollo de nuevos proyectos o la continuación de este mismo si se diera la oportunidad.

\section{Líneas de trabajo futuras}
En esta sección se expondrán las posibles futuras implementaciones que podrían enriquecer el uso de la aplicación así como atraer a nuevos usuarios.

\subsection{Resumen diario}
Se podría incorporar una nueva funcionalidad que informe al usuario sobre su estado en la plataforma, este resumen le brindaría al usuario un resumen diario con información como:
\begin{itemize}
    \item Tareas pendientes para el día actual o los siguientes días.
    \item Nuevas calificaciones en tareas o cuestionarios.
\end{itemize}
Esta funcionalidad nos proporcionaría un enfoque de lo útil que sería usar esta aplicación, ya que el usuario estaría en todo momento informado de sus labores o resultados de las mismas.

\subsection{Visualización de estadísticas}
El hecho de permitir que los usuarios visualicen el progreso en sus tareas o cursos puede servir de ayuda de cara al usuario para concienciarlo en el cumplimiento de sus actividades. Es decir, si un usuario no esta dedicando mucho tiempo a un curso, esta funcionalidad le puede incentivar a ser mas constante y atento. 

Dentro de esta funcionalidad se pueden incluir muchos tipos de estadísticas, estas pueden ser:
\begin{itemize}
    \item Porcentaje de actividades completadas respecto del total.
    \item Tiempo medio dedicado en la realización de actividades.
    \item Media ponderada de las calificaciones de actividades.
\end{itemize}

\subsection{Internacionalización a otro idiomas}
Actualmente la aplicación se encuentra en castellano, por lo que limita su uso a las regiones donde se haga uso de este idioma. El hecho de expandirse por otros horizontes requiere de adaptarse a las necesidades de los usuarios. Es por ello, que uno de los puntos cruciales para poder llegar a más público, es la integración de nuevos idiomas en la aplicación. Como consecuencia se podría comercializar a nivel global.

\subsection{Implementación nuevos diagramas}
Para mejorar la experiencia de gestión de tareas, se propone la integración de nuevos diagramas, gráficos o formas de representación. El hecho de recurrir a nuevas formas de representación de tareas, puede proporcionar al usuario distintas perspectivas para la gestión y planificación de sus labores. Los gráficos a implementar podrían ser:
\begin{itemize}
    \item \textbf{Tablero Kanban:} consiste en un tablero divido en distintas secciones: ``Por hacer``, ``En progreso`` y ``Hecho``, entre otras varias posibilidades. Esta forma de trabajo permite desplazar las tareas entre las distintas secciones dependiendo de su estado.
    \item \textbf{Calendario o agenda:} herramientas que permiten asignar fechas a las tareas con la finalidad de programar plazos o bloques de trabajo.
    \item \textbf{Checklist:} se trata de una lista simple de tareas donde las realizadas quedan marcadas con una casilla, indicando su estado de finalización.
\end{itemize}

\subsection{Sistema de notificaciones}
Un aspecto muy importante a implementar de cara al futuro es la notificación de los eventos a realizar. Es muy importante que el usuario reciba avisos para no perder la constancia en la realización de sus labores. Este sistema se podría plantear de forma que el usuario pueda configurar las notificaciones según sus necesidades o intereses.

\subsection{Implementación de nuevos filtros}
El hecho de que la aplicación se pueda expandir hacia todos los horizontes, hace que el sistema de filtrado cobre mayor importancia. Por ello, la implementación de nuevos filtros podría suponer una interfaz más limpia así como una organización potenciada por dichos filtros. Los filtros que podrían ser implementados son:
\begin{itemize}
    \item Filtro por estado de entrega.
    \item Filtro por prioridad de la tarea (Tareas personales).
    \item Filtro por calificaciones.
    \item Filtro por duración de las tareas.
\end{itemize}

\subsection{Integración con otras plataformas}
Actualmente la aplicación se conecta con las plataformas Moodle, función que permite internacionalizar la aplicación en distintas instituciones que utilicen Moodle. Para llegar a más usuarios, se plantea implementar la compatibilidad con otras plataformas de gestión de tareas y cursos. Algunas plataformas son:
\begin{itemize}
    \item Blackbox
    \item Canvas
    \item Totara
\end{itemize}



\bibliographystyle{plain}
\bibliography{bibliografia}

\end{document}